% GENERATED FILE. Edit canonical lessons, then run npm run generate:books.
% canonical-source-hash: fnv1a64:7e2d68333695385f
% canonical-lessons: RU-C494-suprug, RU-C494-brak, RU-C494-izmena, RU-C494-izmenyat, RU-C494-rasstavaniye

\chapter{Marriage and Break-ups}
\label{ch:ru-marriage-and-break-ups}
\hlchaptermodality{494}

\begin{chapteropening}
\textbf{By the end of this chapter:} \emph{I can talk about marriage, infidelity and breaking up.}

Complete the last lesson of chapter 494: I can talk about marriage, infidelity and breaking up.
\end{chapteropening}

\section[\texorpdfstring{\ru{супруг} (suprúg)}{супруг (suprúg)}]{\ru{супруг} (suprúg) — a spouse, a husband}

\label{lesson:RU-C494-suprug}



\begin{quote}
Before the new one: say the Russian for tenderness, then the Russian for affectionate.
\end{quote}

\subsection*{You'll want to know: \ru{супруг}}
\textbf{\ru{супруг}} (\emph{suprúg}) — \textquotedblleft{}a spouse, a husband\textquotedblright{}.

\begin{quote}
\ru{На} \ru{приём} \ru{она} \ru{пришла} \ru{с} \ru{супругом}.
\end{quote}

\begin{quote}
\emph{Na priyóm oná prishlá s suprúgom.}
\end{quote}

\begin{quote}
She came to the reception with her husband.
\end{quote}

Formal. The wife is \emph{\ru{супруга}}; the plural \emph{\ru{супруги}} means a married couple.

\begin{cousinweb}
\emph{\ru{супруг}} is inherited from Proto-Slavic \emph{sǫprǫg\ru{ъ}} (EtymDB), built on \emph{sǫ-} 'together' and a root usually glossed 'to harness', so spouses are a pair in one harness.
\end{cousinweb}

\begin{grammarlens}[title={What you've built}]
One more word for telling, explaining and describing at B1.
\end{grammarlens}

\subsection*{Your turn}
\begin{itemize}
\raggedright
  \item \emph{Say it:} \emph{suprúg}
  \item \emph{Say it:} \emph{suprúg}, once more
  \item \emph{Say it:} a sentence of your own with \emph{suprúg}
  \item \emph{Recall:} say the Russian for tenderness, then the Russian for affectionate, then say \emph{suprúg} again
\end{itemize}

\begin{tcolorbox}[breakable,colback=teal!4,colframe=teal!35!black,title={Before you move on}]
What does \ru{супруг} mean? (\textquotedblleft{}A spouse, a husband\textquotedblright{}.) Say it once more.
\end{tcolorbox}
\section[\texorpdfstring{\ru{брак} (brak)}{брак (brak)}]{\ru{брак} (brak) — marriage; a defect}

\label{lesson:RU-C494-brak}



\begin{quote}
Before the new one: say the Russian for affectionate, then the Russian for a spouse.
\end{quote}

\subsection*{You'll want to know: \ru{брак}}
\textbf{\ru{брак}} (\emph{brak}) — \textquotedblleft{}marriage; a defect\textquotedblright{}.

\begin{quote}
\ru{Их} \ru{брак} \ru{длится} \ru{уже} \ru{тридцать} \ru{лет}.
\end{quote}

\begin{quote}
\emph{Ikh brak dlítsya uzhé trídtsat let.}
\end{quote}

\begin{quote}
Their marriage has lasted thirty years.
\end{quote}

Masculine. A second word \emph{\ru{брак}} means faulty goods: \emph{\ru{бракованный}} — defective.

\begin{cousinweb}
\emph{\ru{брак}} 'marriage' was borrowed from Church Slavonic \emph{brak\ru{ъ}}, the shape with \emph{-\ru{ра}-} beside Proto-Slavic \emph{bork\ru{ъ}} (EtymDB).
\end{cousinweb}

\begin{grammarlens}[title={What you've built}]
One more word for telling, explaining and describing at B1.
\end{grammarlens}

\subsection*{Your turn}
\begin{itemize}
\raggedright
  \item \emph{Say it:} \emph{brak}
  \item \emph{Say it:} \emph{brak}, once more
  \item \emph{Say it:} a sentence of your own with \emph{brak}
  \item \emph{Recall:} say the Russian for affectionate, then the Russian for a spouse, then say \emph{brak} again
\end{itemize}

\begin{tcolorbox}[breakable,colback=teal!4,colframe=teal!35!black,title={Before you move on}]
What does \ru{брак} mean? (\textquotedblleft{}Marriage; a defect\textquotedblright{}.) Say it once more.
\end{tcolorbox}
\section[\texorpdfstring{\ru{измена} (izména)}{измена (izména)}]{\ru{измена} (izména) — infidelity; betrayal}

\label{lesson:RU-C494-izmena}



\begin{quote}
Before the new one: say the Russian for a spouse, then the Russian for marriage.
\end{quote}

\subsection*{You'll want to know: \ru{измена}}
\textbf{\ru{измена}} (\emph{izména}) — \textquotedblleft{}infidelity; betrayal\textquotedblright{}.

\begin{quote}
\ru{Она} \ru{не} \ru{смогла} \ru{простить} \ru{мужу} \ru{измену}.
\end{quote}

\begin{quote}
\emph{Oná ne smoglá prostít múzhu izménu.}
\end{quote}

\begin{quote}
She couldn't forgive her husband for being unfaithful.
\end{quote}

Feminine. Also treason: \emph{\ru{измена} \ru{родине}}.

\begin{cousinweb}
\emph{\ru{измена}} is built on \emph{\ru{из}-} and the root of \emph{\ru{менять}} 'to change'; EtymDB lists \emph{\ru{измена}} without a source.
\end{cousinweb}

\begin{grammarlens}[title={What you've built}]
One more word for telling, explaining and describing at B1.
\end{grammarlens}

\subsection*{Your turn}
\begin{itemize}
\raggedright
  \item \emph{Say it:} \emph{izména}
  \item \emph{Say it:} \emph{izména}, once more
  \item \emph{Say it:} a sentence of your own with \emph{izména}
  \item \emph{Recall:} say the Russian for a spouse, then the Russian for marriage, then say \emph{izména} again
\end{itemize}

\begin{tcolorbox}[breakable,colback=teal!4,colframe=teal!35!black,title={Before you move on}]
What does \ru{измена} mean? (\textquotedblleft{}Infidelity; betrayal\textquotedblright{}.) Say it once more.
\end{tcolorbox}
\section[\texorpdfstring{\ru{изменять} (izmenyát)}{изменять (izmenyát)}]{\ru{изменять} (izmenyát) — to be unfaithful; to change}

\label{lesson:RU-C494-izmenyat}



\begin{quote}
Before the new one: say the Russian for marriage, then the Russian for infidelity.
\end{quote}

\subsection*{You'll want to know: \ru{изменять}}
\textbf{\ru{изменять}} (\emph{izmenyát}) — \textquotedblleft{}to be unfaithful; to change\textquotedblright{}.

\begin{quote}
\ru{Он} \ru{никогда} \ru{не} \ru{изменял} \ru{жене}.
\end{quote}

\begin{quote}
\emph{On nikogdá ne izmenyál zhené.}
\end{quote}

\begin{quote}
He has never been unfaithful to his wife.
\end{quote}

Imperfective; perfective \ru{изменить}. With the dative, to be unfaithful: \emph{\ru{изменять} \ru{жене}}; with the accusative, to change: \emph{\ru{изменять} \ru{план}}.

\begin{cousinweb}
\emph{\ru{изменять}} is built on \emph{\ru{из}-} and \emph{\ru{менять}} 'to change'; EtymDB lists the verb without a source.
\end{cousinweb}

\begin{grammarlens}[title={What you've built}]
One more word for telling, explaining and describing at B1.
\end{grammarlens}

\subsection*{Your turn}
\begin{itemize}
\raggedright
  \item \emph{Say it:} \emph{izmenyát}
  \item \emph{Say it:} \emph{izmenyát}, once more
  \item \emph{Say it:} a sentence of your own with \emph{izmenyát}
  \item \emph{Recall:} say the Russian for marriage, then the Russian for infidelity, then say \emph{izmenyát} again
\end{itemize}

\begin{tcolorbox}[breakable,colback=teal!4,colframe=teal!35!black,title={Before you move on}]
What does \ru{изменять} mean? (\textquotedblleft{}To be unfaithful; to change\textquotedblright{}.) Say it once more.
\end{tcolorbox}
\section[\texorpdfstring{\ru{расставание} (rasstavániye)}{расставание (rasstavániye)}]{\ru{расставание} (rasstavániye) — a parting, a break-up}

\label{lesson:RU-C494-rasstavaniye}



\begin{quote}
Before the new one: say the Russian for infidelity, then the Russian for to be unfaithful.
\end{quote}

\subsection*{You'll want to know: \ru{расставание}}
\textbf{\ru{расставание}} (\emph{rasstavániye}) — \textquotedblleft{}a parting, a break-up\textquotedblright{}.

\begin{quote}
\ru{После} \ru{расставания} \ru{они} \ru{остались} \ru{друзьями}.
\end{quote}

\begin{quote}
\emph{Pósle rasstavániya oní ostális druzyámi.}
\end{quote}

\begin{quote}
After the break-up they stayed friends.
\end{quote}

Neuter. The verbs are \emph{\ru{расставаться} / \ru{расстаться}} — to part.

\begin{cousinweb}
\emph{\ru{расставание}} is built on \emph{\ru{расставаться}} 'to part', from \emph{\ru{раз}-} 'apart' and \emph{\ru{стать}} 'to stand, to become'; EtymDB has no entry for the noun.
\end{cousinweb}

\begin{grammarlens}[title={What you've built}]
One more word for telling, explaining and describing at B1.
\end{grammarlens}

\subsection*{Your turn}
\begin{itemize}
\raggedright
  \item \emph{Say it:} \emph{rasstavániye}
  \item \emph{Say it:} \emph{rasstavániye}, once more
  \item \emph{Say it:} a sentence of your own with \emph{rasstavániye}
  \item \emph{Recall:} say the Russian for infidelity, then the Russian for to be unfaithful, then say \emph{rasstavániye} again
\end{itemize}

\begin{tcolorbox}[breakable,colback=teal!4,colframe=teal!35!black,title={Before you move on}]
What does \ru{расставание} mean? (\textquotedblleft{}A parting, a break-up\textquotedblright{}.) Say it once more.
\end{tcolorbox}
